\input format.tex

%\usepackage{arydshln}

\begin{document}

\vspace*{3mm}
%%%
\setlength{\arrayrulewidth}{0.5pt}
\fontsize{9.3pt}{17pt}\selectfont
\color{gray2}

\begin{center}
{\bf\sanhao 健康建议}
\end{center}

%\vspace*{1mm}
\begin{spacing}{1.25}
\bigskip
\noindent
根据您本次检测结果，我们从膳食、益生菌、益生元等方面为您制定了肠道菌群干预方案，建议您遵照干预方案进行调理，以促进肠道菌群平衡。并且建议您调理3个月后复检肠道菌群，以便了解干预效果，调整干预方案，帮助您持续提升肠道健康水平。

\bigskip
\tabcolsep=6pt
\fontsize{9.3pt}{9.8pt}\selectfont
\begin{tctabularx}{tabularx={p{15.4cm}<{\centering}}}
\\[-6.769pt]
\cellcolor{topcolor} \raisebox{1pt}{\color{gray0}\xiaosihao\lantxh\bf 膳食方案} \\
\end{tctabularx}

{
\lantxh
\vspace*{-4.25mm}
\noindent
\begin{longtable}{|p{1.25cm}@{}p{2cm}|p{11.732cm}|}
\hline
\parbox[c]{\hsize}{\vskip7pt\raisebox{-1.5ex}{\includegraphics[width=.9cm]{zhushi.pdf}}\vskip7pt} & \parbox[c]{\hsize}{\centering\xiaowuhao \color{gray2}\vskip7pt主食\vskip7pt} & \parbox[c]{\hsize}{\color{gray2}\vskip7pt推荐高膳食纤维饮食。可搭配适量糙米、燕麦、荞麦、玉米及杂豆类等富含膳食纤维的五谷杂粮。\vskip7pt}\\
\hline
\parbox[c]{\hsize}{\vskip7pt\raisebox{-1.5ex}{\includegraphics[width=.9cm]{roulei.pdf}}\vskip7pt} & \parbox[c]{\hsize}{\centering\xiaowuhao \color{gray2}\vskip7pt肉类\vskip7pt} & \parbox[c]{\hsize}{\color{gray2}\vskip7pt宜选择瘦猪肉、牛肉、去皮禽肉，避免高脂肪高胆固醇的肥猪肉、肥牛、动物内脏等。\vskip7pt}\\
\hline
\parbox[c]{\hsize}{\vskip7pt\raisebox{-1.5ex}{\includegraphics[width=.9cm]{yuxialei.pdf}}\vskip7pt} & \parbox[c]{\hsize}{\centering\xiaowuhao \color{gray2}\vskip7pt鱼虾类\vskip7pt} & \parbox[c]{\hsize}{\color{gray2}\vskip7pt适量吃鱼虾类优质蛋白，多采用清蒸、水煮的烹饪方式，如清蒸鱼、白灼虾、水煮鱼片等。\vskip7pt}\\
\hline
\parbox[c]{\hsize}{\vskip7pt\raisebox{-1.5ex}{\includegraphics[width=.9cm]{danlei.pdf}}\vskip7pt} & \parbox[c]{\hsize}{\centering\xiaowuhao \color{gray2}\vskip7pt蛋奶、大豆类\vskip7pt} & \parbox[c]{\hsize}{\color{gray2}\vskip7pt每天一个鸡蛋，奶类宜选用酸奶，适量摄入黄豆、豌豆、蚕豆、豆腐、豆浆等豆类及豆制品。\vskip7pt}\\
\hline
\parbox[c]{\hsize}{\vskip7pt\raisebox{-1.5ex}{\includegraphics[width=.9cm]{shucaishuiguo.pdf}}\vskip7pt} & \parbox[c]{\hsize}{\centering\xiaowuhao \color{gray2}\vskip7pt蔬菜水果\vskip7pt} & \parbox[c]{\hsize}{\color{gray2}\vskip7pt多吃富含膳食纤维的新鲜果蔬，如柚子、柑橘、苹果、香蕉、火龙果、洋葱、芹菜、白菜等。\vskip7pt}\\
\hline
\parbox[c]{\hsize}{\vskip7pt\raisebox{-1.5ex}{\includegraphics[width=.9cm]{you.pdf}}\vskip7pt} & \parbox[c]{\hsize}{\centering\xiaowuhao \color{gray2}\vskip7pt食用油、盐\vskip7pt} & \parbox[c]{\hsize}{\color{gray2}\vskip7pt以橄榄油、菜籽油作为膳食油脂的主要来源，食用油每天不超过25g，食盐摄入量每天不超过6g。\vskip7pt}\\
\hline
\parbox[c]{\hsize}{\vskip7pt\raisebox{-1.5ex}{\includegraphics[width=.9cm]{other.pdf}}\vskip7pt} & \parbox[c]{\hsize}{\centering\xiaowuhao \color{gray2}\vskip7pt其他\vskip7pt} & \parbox[c]{\hsize}{\color{gray2}\vskip7pt蔬菜的烹饪方式以轻烫为宜，不宜过度加工。如患痛风、溃疡性结肠炎、慢性腹泻、食物不耐受等有饮食禁忌的疾病，请优先遵循疾病的饮食原则。\vskip7pt}\\
\hline
\end{longtable}
\vspace*{-6mm}
\begin{longtable}{p{0.1cm}p{15.1cm}}
\includegraphics{xiaofangkuai.pdf} & {\qihao 以上膳食方案仅供参考，您的实际饮食请结合自身具体情况调整。} \\
\end{longtable}
}
\end{spacing}
\end{document}
